\documentclass[a4paper]{article}

% --- Packages ---
\usepackage{amsmath, amssymb} % Equations and math symbols
\usepackage{graphicx}         % Figures
\usepackage{listings}         % Code blocks
\usepackage{cite}             % BibTeX citations
\usepackage{hyperref}         % Clickable references

% --- Code Block Configuration ---
\lstset{
    basicstyle=\ttfamily\footnotesize,
    breaklines=true,
    numbers=left,
    numberstyle=\tiny,
    frame=single
}

\title{Minimal LaTeX Template}
\author{Author Name}
\date{\today}

\begin{document}

\maketitle

\begin{abstract}
This is a minimal boilerplate for an academic article demonstrating the integration of equations, figures, code snippets, and BibTeX citations.
\end{abstract}

\section{Introduction}
This template relies on standard packages. An example citation is placed here \cite{lamport1994}.

\section{Equations}
The governing equation is presented in Equation \ref{eq:euler}:
\begin{equation} \label{eq:euler}
    e^{i\pi} + 1 = 0
\end{equation}

\section{Figures}
Figure \ref{fig:placeholder} demonstrates the figure environment.
\begin{figure}[htbp]
    \centering
    % Uncomment and replace 'filename' with your image file
    % \includegraphics[width=0.5\textwidth]{filename}
    \rule{5cm}{3cm} % Black box placeholder
    \caption{Placeholder for an embedded graphic.}
    \label{fig:placeholder}
\end{figure}

\section{Code Listings}
A standard code implementation is shown in Listing \ref{lst:example}.

\begin{lstlisting}[language=Python, caption={Python example}, label={lst:example}]
import sys

def compute_factorial(n: int) -> int:
    if n == 0:
        return 1
    return n * compute_factorial(n - 1)

print(compute_factorial(5))
\end{lstlisting}

% --- Bibliography ---
\bibliographystyle{plain}
\bibliography{references}

\end{document}
